% Apêndice do guia: listagens de código com o pacote listings. Apague-o no seu trabalho.

% definição da linguagem json, usada nos exemplos (arquivo exibido na Seção D.4)
\lstdefinelanguage{json}{
  morestring=[b]",
  morecomment=[l]{//},
  morekeywords={true, false, null},
  sensitive=true}


\chapter{Guia de listagens de código}\label{ap:listagens}

Este apêndice reúne a configuração e as instruções básicas do pacote \texttt{listings},
usado para apresentar código-fonte no trabalho. A configuração do guia (no preâmbulo do
arquivo \texttt{abntex2-ufv.tex}; uma cópia comentada está em \texttt{template/template.tex})
define uma fonte de largura fixa legível, a numeração das linhas e a sintaxe colorida. As
listagens são ilustrações do tipo ``Código'': ficam dentro dos ambientes \texttt{codigo} e
\texttt{ilustracao}, com título, fonte e, se preciso, legenda e nota
(Apêndice~\ref{ap:tabelas}).

\section{Configuração recomendada}\label{sec:lst-config}

O trecho abaixo, do preâmbulo, escolhe a fonte, define as cores e cria o estilo
\texttt{ufvcodigo}, aplicado a todas as listagens com \texttt{\textbackslash lstset}:

\begin{lstlisting}[style=exemplo, language={[LaTeX]TeX}]
\usepackage[scaled=0.88,lining]{FiraMono}   % fonte de largura fixa legível
\usepackage{xcolor}
\usepackage{listings}
\definecolor{ufvkw}{RGB}{0,51,153}      % palavras-chave
\definecolor{ufvcom}{RGB}{0,110,50}     % comentários
\definecolor{ufvstr}{RGB}{163,21,21}    % cadeias de caracteres
\lstdefinestyle{ufvcodigo}{
  basicstyle=\ttfamily\footnotesize, columns=fixed, keepspaces=true, tabsize=4,
  keywordstyle=\bfseries\color{ufvkw}, commentstyle=\itshape\color{ufvcom},
  stringstyle=\color{ufvstr}, showstringspaces=false,
  numbers=left, numberstyle=\tiny\color{black!55}, numbersep=7pt,
  backgroundcolor=\color{black!4}, frame=single, rulecolor=\color{black!25},
  breaklines=true, extendedchars=true,
  literate={á}{{\'a}}1 {é}{{\'e}}1 {ã}{{\~a}}1 {ç}{{\c c}}1 ...}   % acentos
\lstset{style=ufvcodigo}
\end{lstlisting}

\begin{itemize}
  \item \textbf{Fonte}: \texttt{FiraMono} distingue bem 0/O e 1/l/I, mantém a largura fixa e
        tem negrito, itálico e negrito-itálico de verdade, no \texttt{pdflatex}, no
        Lua\LaTeX\ e no Xe\LaTeX. A opção \texttt{lining} dá algarismos alinhados (sem ela, a
        fonte usa algarismos de estilo antigo, ruins em código) e \texttt{scaled=0.88} casa a
        altura das letras com a do Times e deixa cerca de 85 colunas por linha no
        \texttt{\textbackslash footnotesize}. Alternativas com negrito e itálico reais, também
        instaladas no MiKTeX e no TeX Live: \texttt{beramono} com \texttt{scaled=0.81} (casa a
        altura das letras com a do Times e deixa cerca de 93 colunas, com letras um pouco mais
        miúdas; com \texttt{scaled=0.9} as letras ficam altas demais e cabem só 84 colunas),
        \texttt{sourcecodepro} (\texttt{scaled=0.88}) e \texttt{plex-mono}
        (\texttt{scale=0.88}). Evite o \texttt{inconsolata} (\texttt{zi4}) e o
        \texttt{lmodern}, sem itálico e sem negrito de largura fixa, respectivamente (o \LaTeX\
        acusa \emph{Font shape undefined}). Como o pacote muda o
        \texttt{\textbackslash ttdefault}, vale também para o \texttt{\textbackslash texttt}
        do texto.
  \item \textbf{Largura fixa}: use \texttt{columns=fixed} (e não \texttt{flexible}), senão o
        alinhamento das colunas do código se perde.
  \item \textbf{Cores}: as ilustrações podem ser coloridas (manual UFV 4.1), mas escolha
        cores que continuem legíveis em escala de cinza; palavras-chave em negrito e
        comentários em itálico ajudam. Para quem imprime em preto e branco, o estilo
        \texttt{ufvcodigopb} (Seção~\ref{sec:cod-pb}) desliga as cores.
\end{itemize}

\section{Uso básico}

O código vai dentro do ambiente \texttt{lstlisting}, com a linguagem em
\texttt{[language=...]}. O resultado do \autoref{cod:guia-python} usa o estilo padrão:

\codexemplo{codigo-python}

Para listar parte de um arquivo, em vez de copiar o código, use \texttt{\textbackslash lstinputlisting}
com \texttt{firstline} e \texttt{lastline} (ou \texttt{linerange}); assim o código do texto
nunca fica diferente do arquivo (\autoref{cod:guia-arquivo}):

\codexemplo{codigo-arquivo}

Para um trecho curto no meio de uma frase, use \texttt{\textbackslash lstinline}: por
exemplo, \lstinline[language=Python]{total += v} soma o valor ao acumulador. Como o código
vem entre delimitadores (aqui, chaves), evite o delimitador dentro do trecho.

\section{Opções mais usadas}

O \autoref{qua:lst-opcoes} resume as opções mais úteis. Elas podem ser dadas no
\texttt{\textbackslash lstset} (para todas as listagens) ou entre colchetes em uma listagem
(só para ela).

\begin{quadro}[htbp]
\begin{ilustracao}[\linewidth]{\caption[Opções mais usadas do pacote listings]{Opções mais usadas do pacote \texttt{listings}}\label{qua:lst-opcoes}}{\elaboracaopropria}
  \renewcommand{\arraystretch}{1.2}
  \small
  \begin{tabularx}{\linewidth}{|>{\ttfamily\raggedright\arraybackslash}p{4.2cm}|X|}
    \hline
    \textnormal{\textbf{Opção}} & \textbf{Efeito} \\
    \hline
    language=Python & Linguagem da listagem (veja a Seção~\ref{sec:lst-linguagens}) \\
    \hline
    numbers=none & Remove a numeração das linhas (\texttt{left} a coloca à esquerda) \\
    \hline
    firstnumber=10 & Começa a numeração na linha 10 (útil para trechos de arquivos) \\
    \hline
    firstline=4, lastline=11 & Em \texttt{\textbackslash lstinputlisting}, mostra só as linhas 4 a 11 \\
    \hline
    frame=single & Moldura (\texttt{none}, \texttt{tb}, \texttt{lines} etc.) \\
    \hline
    backgroundcolor & Cor de fundo da listagem, por exemplo \texttt{\textbackslash color\{black!4\}} \\
    \hline
    breaklines=true & Quebra linhas longas (e \texttt{postbreak} marca a continuação) \\
    \hline
    basicstyle & Fonte e tamanho do código, por exemplo \texttt{\textbackslash ttfamily\textbackslash footnotesize} (tamanho 10); \texttt{\textbackslash small} dá letra um pouco maior \\
    \hline
    keywordstyle, commentstyle, stringstyle & Estilo de palavras-chave, comentários e cadeias (cor, negrito, itálico) \\
    \hline
    morekeywords=\{foo,\allowbreak bar\} & Acrescenta palavras-chave à linguagem \\
    \hline
    escapechar=\% & Permite comandos \LaTeX\ dentro do código, depois do caractere (por exemplo, \texttt{\textbackslash textbf}) \\
    \hline
    tabsize=4 & Largura do caractere de tabulação, em espaços \\
    \hline
    linerange=\{3-5,9-12\} & Em \texttt{\textbackslash lstinputlisting}, mostra intervalos de linhas \\
    \hline
  \end{tabularx}
\end{ilustracao}
\end{quadro}

\section{Linguagens}\label{sec:lst-linguagens}

O \texttt{listings} reconhece, entre outras, \texttt{Python}, \texttt{SQL}, \texttt{Java},
\texttt{C}, \texttt{C++}, \texttt{HTML}, \texttt{XML}, \texttt{PHP}, \texttt{Bash}
(\texttt{bash}), \texttt{Matlab} e \texttt{TeX}. Para uma linguagem que o pacote não conhece,
como \texttt{JSON}, defina-a uma vez com \texttt{\textbackslash lstdefinelanguage}, listando
as palavras-chave, os comentários e as cadeias:

\lstinputlisting[style=exemplo, language={[LaTeX]TeX}]{example/snippets/lang-json.tex}

e use-a como qualquer outra (\autoref{cod:guia-json}):

\codexemplo{codigo-json}

Se a linguagem tiver variantes (dialetos), escolha-a com a sintaxe
\texttt{[language=[ANSI]C]}. A lista completa está no manual do pacote (\texttt{texdoc listings}).

\section{Acentos e caracteres especiais}

Com o arquivo em UTF-8 e \texttt{pdflatex}, o \texttt{listings} não reconhece os
caracteres acentuados por si só; a opção \texttt{literate}, na configuração da
Seção~\ref{sec:lst-config}, converte cada letra acentuada na sua versão correta (é preciso listar as letras
que aparecem no código, em maiúsculas e minúsculas). Sem ela, o PDF mostra caracteres
estranhos ou o \LaTeX\ acusa erro (\emph{Invalid UTF-8 byte sequence}). Com Lua\LaTeX\ ou
Xe\LaTeX\ a opção é desnecessária, mas inofensiva: pode ficar na configuração. Se o código tiver símbolos como \texttt{<}, \texttt{>} ou
\texttt{|} que saiam errados em algumas fontes, use fonte \texttt{T1} (\texttt{\textbackslash usepackage[T1]\{fontenc\}}),
já carregada neste modelo.

\section{Cores, impressão em preto e branco e alternativas}\label{sec:cod-pb}

O estilo \texttt{ufvcodigopb} é o mesmo \texttt{ufvcodigo} sem cores (palavras-chave em
negrito, comentários em itálico, fundo branco). Escolha-o em uma listagem com
\texttt{[style=ufvcodigopb]} ou para todas com \texttt{\textbackslash lstset\{style=ufvcodigopb\}}
(\autoref{cod:guia-pb}):

\codexemplo{codigo-pb}

Outra opção é o pacote \texttt{minted}, que usa o \emph{Pygments} (Python) e produz uma
coloração mais fiel à de editores de código, para muitas linguagens. Ele exige o Python
instalado e a compilação com \texttt{-shell-escape}; o Overleaf suporta, mas em máquinas
locais isso complica a instalação e a segurança. O \texttt{listings} dispensa tudo isso e é
a escolha deste modelo.

\section{Dicas}

\begin{itemize}
  \item Código longo (mais de uma página) vai para apêndice ou repositório; no texto, mostre
        só o essencial e cite o repositório nas referências.
  \item Nunca use imagem (\emph{print}) de código: não pode ser copiado nem lido por leitores
        de tela, e perde nitidez.
  \item Em \texttt{ilustracao}, use o argumento opcional \texttt{[\textbackslash linewidth]}: o
        conteúdo é composto nessa largura, então as linhas longas são quebradas
        (\texttt{breaklines=true}) em vez de estourar a margem, e o título e o rodapé
        ficam alinhados à moldura da listagem. Não é preciso envolver a listagem em
        \texttt{minipage} (documentos que já o fazem continuam funcionando).
  \item Termos longos no texto, como parâmetros de uma API, não quebram em \texttt{\textbackslash texttt}
        (a fonte de largura fixa desliga a hifenização) e podem estourar a margem. Escreva-os
        com \texttt{\textbackslash texttt\{\textbackslash path\{/api/v1/projetos/12345/tarefas\}\}}:
        o comando \texttt{\textbackslash path} (do pacote \texttt{url}, já carregado pelo
        \texttt{hyperref}) permite a quebra de linha, e com o \texttt{xurl} do modelo a
        quebra pode ocorrer entre quaisquer letras. Para só evitar o link, use
        \texttt{\textbackslash nolinkurl\{...\}}; em último caso, divida o trecho em dois
        \texttt{\textbackslash texttt}.
  \item Cuide da indentação e da largura: com \texttt{\textbackslash footnotesize} cabem cerca de
        85 colunas de código na largura do texto; linhas mais longas são quebradas
        (\texttt{breaklines=true}).
  \item Sempre cite a listagem no texto, com \texttt{\textbackslash autoref}, e explique o que
        ela faz.
\end{itemize}
